2月の個別releaseを一列に並べるだけでは、この月の意味は見えにくい。Codex appは複数coding agentとlong-running workを管理するworkspaceを提示し、Qwen3-Coder-Nextはmulti-turn tool interactionとrepository-level understandingを対象にし、Voxtral Realtimeはstreaming speech、Deep Thinkはscience/research/engineering向けspecialized reasoning、vLLMはRealtime APIやResponses/tool surfaceをruntimeへ取り込んだ。対象は異なるが、共通してmodelの出力品質だけでなく『どの時間幅・tool・modality・runtimeで仕事として成立させるか』へ技術課題が広がっている。これが2月全体を貫く転換点である。\autocite{src-15823ae5556e,src-5041feb6f87f,src-d8b7c1c36830,src-dd5d9b8662f3,src-eac8dfbeb78c}


2月から後続月へ持ち越される観察軸は五つある。

\begin{itemize}
  \item agent execution：一回の回答ではなく、複数agent・tool use・long-running taskを安全に継続できるか。
  \item serving/runtime：model capabilityがRealtime API、tool protocol、scheduler、hardware supportまで到達するか。
  \item safety/control：agentが長く強く動くほど、実行権限とcyber boundaryをどう設計するか。
  \item multimodal interaction：speechやimageが低latency・高頻度のinteraction loopへ定着するか。
  \item specialized reasoning：汎用modelとは別に、science/research/engineeringなど特定workloadへreasoningを特化する設計がどう位置付くか。
\end{itemize}

これらは3月以降の具体的イベントを2月へ持ち込むためではなく、2月時点で成立した技術的な問いを後続月でも同じ物差しで読むための軸である。
\autocite{src-15823ae5556e,src-4a841ea7ec7c,src-5041feb6f87f,src-9aba7995b6a5,src-d8b7c1c36830,src-dd5d9b8662f3,src-dd64fe3084a4}

